\chapter{Evaluation}
\label{ch:evaluation}

This chapter evaluates the decomposition of Chapter~\ref{ch:ivm}, realised as the SQL pipelines of Chapter~\ref{ch:implementation}, against a flat-join baseline.
The baseline is not an unmaintained query but the same query handed to Feldera as a single flat view, which the engine keeps up to date incrementally exactly as it does the decomposition's bag views.
The engine and the data are held fixed, so what we isolate is the effect of the rewrite alone: whether maintaining the query as its decomposition's views is cheaper, in throughput and peak memory, than maintaining it as one flat view. The variants that materialise the full output all return the same result, which doubles as a correctness check (Section~\ref{sec:methodology}).

The reported metrics and the pipeline notation are defined in full here, expanding the brief description given with the harness (Section~\ref{sec:harness}). The experimental setup, datasets, hypotheses, and results then follow in turn.

\section{Experimental Setup}
\label{sec:setup}

All experiments were run on \texttt{dast-cpu05}, a CPU-only node of the DaST research cluster at the Department of Informatics (IfI), University of Zurich.

The node has two Intel Xeon Silver 4214 processors at $2.20$\,GHz ($12$ cores each, for $24$ physical cores and $48$ hardware threads), $196$\,GB of main memory and a local SSD. Datasets reside on the node's local scratch storage ($492$\,GB) and are read locally, so that data access introduces no network overhead. The host runs Debian GNU/Linux 12 (bookworm) with Linux kernel \texttt{6.1.0-45-amd64}.

The Feldera engine was deployed on the same host from the official Feldera Docker image (\texttt{pipeline-manager}) in its default configuration, with no changes to the engine itself other than the volume mounts described in Section~\ref{sec:harness}: the container's storage directory is bind-mounted to host scratch to hold the large intermediate views, and the dataset directory is mounted read-only into the container. We pin Feldera to version $0.286.0$ in both \texttt{docker-compose.yaml} (the engine image) and \texttt{requirements.txt} (the Python client), so that the engine and the harness driving it move together. The host ran Docker Engine $29.7.0$ and Python $3.11$. The harness runs directly on the host and drives the containerised engine through its REST API. Every variant begins from a cold start, with the engine's storage wiped between variants (Section~\ref{sec:harness}), so that no residual state or cache from an earlier run can influence a measurement.

The datasets were streamed from local CSV files into Feldera through the \texttt{file\_input} connectors described in Section~\ref{sec:pipeline}, simulating a high-velocity, insert-only ingestion.

\section{Metrics and Pipelines}
\label{sec:methodology}

\subsection{Metrics}

As described in Section~\ref{sec:harness}, the harness samples the running pipeline every five seconds, and once more when it finishes, recording at each sample its uptime, the number of input records processed so far, and its resident-set size (RSS).
Every figure we report is derived from this polling.
From a finished run we take three metrics: the \emph{time}, the wall-clock duration until the pipeline has ingested the whole dataset and its views are up to date, the \emph{throughput}, the number of input records processed per second, computed as the total records over that uptime, and the \emph{peak memory}, the largest RSS seen across all samples of the run.
Alongside these we report the size of the materialised \texttt{RESULT} view, in rows, and the \emph{speedup}, the ratio of the baseline's time to the variant's time, which expresses how many times faster (or, when below one, slower) a variant maintains the query than the flat baseline.

Every experiment is run as three independent trials, with the engine's storage cleared between them so that each trial starts from a clean state.
We report the time and throughput as the mean over the three trials, and the peak memory as the maximum over them, so that the memory figure reflects the worst case a run reached rather than a smoothed average.
An entry marked \texttt{TO} denotes an experiment in which all three trials hit the three-hour timeout without finishing.
For such an experiment we still report the peak memory the engine reached and, in parentheses, the mean fraction of the input processed when the timeout struck, but not its output size, throughput, or speedup, since the run never completed. A run that instead crashes after exhausting the storage volume is marked \texttt{OOS} and reported the same way.
The baseline is the reference for the speedup and so is always $1.00\times$. When the baseline itself times out, its true time exceeds the three-hour timeout, so a variant that finished is reported with a lower-bound speedup $>\!N\times$, where $N$ is the timeout divided by the variant's time.
Table~\ref{tab:metrics} summarises the metrics and how each is aggregated.

\begin{table}[H]
  \centering
  \begin{tabular}{@{}lll@{}}
    \toprule
    \textbf{Metric} & \textbf{Per trial} & \textbf{Reported} \\
    \midrule
    Time        & final uptime                          & mean over the 3 trials \\
    Throughput  & records processed $\div$ final uptime  & mean over the 3 trials \\
    Peak memory & largest RSS among the trial's samples  & maximum over the 3 trials \\
    Output size & rows in the \texttt{RESULT} view       & identical across trials \\
    Speedup     & ---                                    & baseline time $\div$ variant time \\
    \bottomrule
  \end{tabular}
  \caption{The reported metrics and how each is aggregated over the three trials. On a run that did not finish, whether timed out (\texttt{TO}) or out of storage (\texttt{OOS}), only peak memory is reported.}
  \label{tab:metrics}
\end{table}

\medskip
\label{sec:caveats}
A few measurement artefacts should be kept in mind when reading the tables:
\begin{itemize}
    \item \textbf{Averaging noise.} The reported times are the mean of three trials, so differences on the order of $0.05$\,min lie within trial-to-trial variance and should not be over-interpreted.
    \item \textbf{Memory sampling.} Peak memory is sampled every five seconds, so for a very short run (for example $0.10$\,min) the reported value may reflect a single, possibly unrepresentative, snapshot rather than the true peak.
\end{itemize}

\subsection{Pipelines and Variants}
\label{sec:variants}

Every query is run as a flat baseline and as a family of decomposition variants. In the result tables (Tables~\ref{tab:results_soc-Epinions1}--\ref{tab:results_imdb_job}) each decomposition variant is named by a compact symbol recording three choices: how the decomposition tree is rooted, whether it uses the standard IVM$^+$ decomposition or our variant, and how much of the result the pipeline reconstructs.

\begin{itemize}
    \item \textbf{Baseline}: the query as a single flat join, with no decomposition. This is the reference against which the decomposition variants are compared.
    \item \textbf{Orientation ($D_1$, $D_2$)}: the bag relations and their $P'$-semijoin filters are defined over a \emph{rooted} tree decomposition (Section~\ref{sec:bag_relations}), but the decomposition itself does not fix the root. Re-rooting changes which bag is a parent of which, and so the direction in which the $P'$-relations filter and the order in which the bags are maintained. $D_1$ and $D_2$ denote two orientations of the same decomposition: $D_1$ is the less balanced one, a deeper, more chain-like tree with the larger relations in the lower bags, and $D_2$ is the more balanced one, with the bags spread across independent branches and the larger relations nearer the root. Where a decomposition admits only one meaningful orientation, for example a two-bag one, we report only $D_1$.
    \item \textbf{IVM$^+$ ($D$) or the variant ($D'$)}: the unprimed $D$ builds each bag query the standard IVM$^+$ way, keeping every atom that shares a variable with the bag and projecting the partially overlapping ones onto the bag's variables as semijoin filters (Section~\ref{sec:ivm_algorithm}). The primed $D'$ is our variant, which drops those projected atoms and keeps only the atoms fully contained in the bag (Section~\ref{sec:variant_no_semijoin}). The two give the same result and differ only in intermediate size and join work.
    \item \textbf{Reconstruction ($D^{\text{full}}$, $D$, $D^{\text{count}}$)}: $D^{\text{full}}$ performs the reconstruction join over all the bags to materialise the complete result (Section~\ref{sec:bag_relations}). The plain $D$, shown indented beneath its $D^{\text{full}}$, stops at the root bag and skips the reconstruction, maintaining only the bag hierarchy. $D^{\text{count}}$ returns the size of the full result through the counting semiring, without ever materialising the output (Section~\ref{sec:semiring}).
\end{itemize}

After every trial the harness verifies the run by reading the row count of the maintained output view (Section~\ref{sec:harness}). For the baseline and the full-reconstruction variants this is the size of the complete result, which they must all agree on. The without-reconstruction variants materialise only the root bag, so their reported count is that bag's size rather than the full result, and the count variant is checked against the full result size, which it must reproduce through the counting semiring without materialising the output.

Because the baseline maintains the full join result as a single materialised view, a like-for-like comparison must ask the decomposition to produce that same output. The full-reconstruction variant ($D^{\text{full}}$) does so, re-joining the bags into the complete result, so it and the baseline maintain the same answer and differ only in how they compute it. The without-reconstruction variant ($D$) instead stops at the root bag and never forms the full output. This is closer to how incremental maintenance is normally used, keeping a compact structure the result can be read from on demand rather than the result itself (Chapter~\ref{ch:introduction}), so it is not a like-for-like competitor to the baseline but a measure of how cheaply the decomposition alone can be maintained. The count variant ($D^{\text{count}}$) lies between the two, reporting the result's size without materialising it.

\medskip
All experiments were run as pipelines in which every table and every view is materialised, so that any of them can be queried directly. In practice this is often unnecessary, so it is a conservative choice, and one that weighs more heavily on the decomposition, which introduces many additional views. As described in Section~\ref{sec:feldera}, materialising a table or view can slow a pipeline down, so any advantage the decomposition shows here is obtained despite this. To gauge the effect, we ran a subset of the experiments in which only each pipeline's terminal view is materialised (Section~\ref{sec:pipeline}). Tables with a primary key stay materialised regardless, since Feldera maintains those automatically (Section~\ref{sec:feldera}). These runs are reported separately, one table per dataset, in Appendix~\ref{app:results_ablation}.

\section{Datasets and Queries}
\label{sec:datasets}

We evaluate on three classes of dataset, introduced here together with the queries run on each.

\subsection{Graph datasets}
The graph datasets are three directed graphs from the SNAP collection~\cite{snapnets}, each stored as a single \texttt{EDGES} relation. \texttt{soc-Epinions1} is a who-trusts-whom social network ($|V| = 75'879$, $|E| = 508'837$), \texttt{cit-HepTh} an arXiv high-energy-physics citation network ($|V| = 27'770$, $|E| = 352'807$), and \texttt{p2p-Gnutella31} a Gnutella peer-to-peer snapshot ($|V| = 62'586$, $|E| = 147'892$). Every query is a self-join over the edge relation: one acyclic query and several cyclic queries (Table~\ref{tab:queries_graph}).

Because every atom is the same \texttt{EDGES} relation, the projections that separate IVM$^+$ from the variant only filter a bag down to the nodes with an edge on one endpoint, almost every node in a dense graph, and so barely filter while still costing extra work. We therefore run only the variant on these datasets ($D'$).

Each query is run under two input configurations. In the first, the edge file is ingested through a single connector into one \texttt{EDGES} table, which the query aliases once per atom. In the second, each atom has its own table, all loaded with the same edge file, so the query remains a self-join but over separate copies of the data. Both configurations produce the same result and are reported separately, labelled by the number of connectors. We compare them on time rather than throughput, since the multi-connector runs re-ingest the same file and inflate the record count.

\begin{table}[H]
  \centering
  \begin{tabular}{@{}llc@{}}
    \toprule
    \textbf{Query} & \textbf{Structure} & $\boldsymbol{w(Q)}$ \\
    \midrule
    5 Path Query   & path, five edges          & $1$ \\
    Cyclic Query 1 & 6-cycle with three chords & $\tfrac{3}{2}$ \\
    Cyclic Query 2 & 6-cycle with two chords   & $2$ \\
    Cyclic Query 3 & diamond ($Q_\diamond$)    & $\tfrac{3}{2}$ \\
    Cyclic Query 4 & 6-cycle                   & $2$ \\
    Cyclic Query 5 & 5-cycle                   & $2$ \\
    Cyclic Query 6 & 4-cycle                   & $2$ \\
    \bottomrule
  \end{tabular}
  \caption{The queries run on the graph datasets, with their fractional hypertree width $w(Q)$. Full definitions are given in Appendix~\ref{app:queries}. Cyclic Query 3 is the running example $Q_\diamond$ (Section~\ref{sec:queries}).}
  \label{tab:queries_graph}
\end{table}

\subsection{TPC-H}
TPC-H~\cite{poess2000} is a standard decision-support benchmark over an eight-relation schema. We generate it at three scale factors, 1, 10, and 30, and run three acyclic queries and one cyclic query (Table~\ref{tab:queries_tpch}).

\begin{table}[H]
  \centering
  \begin{tabular}{@{}llc@{}}
    \toprule
    \textbf{Query} & \textbf{Structure} & $\boldsymbol{w(Q)}$ \\
    \midrule
    3 Table Path Query & path over three tables & $1$ \\
    4 Table Path Query & path over four tables  & $1$ \\
    5 Table Path Query & path over five tables  & $1$ \\
    Cyclic Query 7     & 4-cycle over four tables & $2$ \\
    \bottomrule
  \end{tabular}
  \caption{The queries run on TPC-H, with their fractional hypertree width $w(Q)$. Full definitions are given in Appendix~\ref{app:queries}.}
  \label{tab:queries_tpch}
\end{table}

\subsection{IMDB (JOB)}
The Join Order Benchmark (JOB)~\cite{leis2015} is derived from the real Internet Movie Database, and unlike the synthetic TPC-H data it contains real-world, correlated data. We run two acyclic queries and one cyclic query (Table~\ref{tab:queries_imdb}).

\begin{table}[H]
  \centering
  \begin{tabular}{@{}llc@{}}
    \toprule
    \textbf{Query} & \textbf{Structure} & $\boldsymbol{w(Q)}$ \\
    \midrule
    Acyclic Query 1 & path over five tables      & $1$ \\
    Acyclic Query 2 & tree over ten tables       & $1$ \\
    Cyclic Query 8  & 3-cycle over three tables & $\tfrac{3}{2}$ \\
    \bottomrule
  \end{tabular}
  \caption{The queries run on IMDB (JOB), with their fractional hypertree width $w(Q)$. Full definitions are given in Appendix~\ref{app:queries}.}
  \label{tab:queries_imdb}
\end{table}

\section{Hypotheses}
\label{sec:hypotheses}

The experiments were designed around five hypotheses derived from the theory of Chapters~\ref{ch:background}--\ref{ch:ivm}.

\begin{description}
    \item[\hypertarget{h1}{H1 (Tree orientation).}] A tree decomposition fixes the bags but not the root. Re-rooting changes the parent--child relationships between bags and therefore the direction in which the $P'$-semijoin filters propagate and the order in which the bags are computed bottom-up. A child propagates $P'_c = \pi_{Y_c} Q'_c$ to its parent; the smaller (more selective) this $P'$ relation, the more parent tuples the semijoin eliminates and the cheaper it is to maintain. We therefore expect the orientation that lets small, selective filters flow toward a heavy root, reducing the size of the largest materialised bag relation, to perform best and \emph{a priori} expect the more balanced orientation (\textbf{D2}) to achieve this better than the deeper, more chain-like orientation (\textbf{D1}).

    \item[\hypertarget{h2}{H2 (Input-connector strategy).}] For the self-join graph queries, ingesting the edge relation once through a single connector and aliasing it in SQL should perform comparably to declaring one connector per relation, because the dominant cost is tuple \emph{processing} rather than ingestion. The single-connector variant may be marginally cheaper as the data is ingested and parsed only once.

    \item[\hypertarget{h3}{H3 (Acyclic queries).}] Acyclic queries have $w(Q)=1$, so an optimal decomposition places one relation per bag, closely mirroring the join tree a query optimiser would build for the baseline. We therefore expect the \emph{full-reconstruction} decomposition ($D^{\text{full}}$) to offer no throughput advantage over the baseline while paying a memory overhead, as it materialises every bag view in addition to the final result. The \emph{without-reconstruction} decomposition ($D$) should instead be cheap, since it only maintains the bag hierarchy and exposes the root bag. We expect the reconstruction step to dominate the cost of $D^{\text{full}}$. Beyond the unavoidable cost of materialising a potentially large output, the final re-join of all bag views is executed by Feldera as an ordinary SQL join, with the join order and algorithm chosen by its optimiser; since the engine compiles joins to binary (pairwise) operators rather than a worst-case-optimal multiway join such as leapfrog triejoin, and we do not control the order in which it re-joins the bags, this reconstruction may be carried out far less efficiently than the decomposition's tree structure would in principle allow. We further expect the inline semijoins of the \emph{with-semijoins} variant to add redundant maintenance work unless they act as strong filters.

    \item[\hypertarget{h4}{H4 (Cyclic queries).}] Cyclic queries have $w(Q)>1$, and are precisely the queries for which a flat binary-join plan can produce large intermediate results. We expect the baseline to suffer exactly this intermediate explosion, whereas the decomposition avoids it: rather than materialising the full flat join, it groups relations into bags and connects them through the projected $P'$-relations, bounding the largest intermediate to roughly the heaviest bag's local join instead of the flat plan's blow-up.\footnote{The clean $N^{w(Q)-1}$ bound assumes each bag is evaluated optimally. Since Feldera also uses binary joins \emph{within} a bag, the realised benefit comes from the bags being smaller sub-patterns connected by cheap projections rather than from a guaranteed worst-case-optimal bound.}
    The variant-level overheads are the same as in \hyperlink{h3}{H3}: the \emph{without-reconstruction} variant ($D$) is cheap, the \emph{with-semijoins} variant adds redundant maintenance work unless its projections filter strongly, and the \emph{full-reconstruction} variant ($D^{\text{full}}$) pays the reconstruction cost discussed there. Unlike the acyclic case, however, we expect the intermediate-bounding benefit to outweigh these overheads: the decomposition should beat the baseline even under full reconstruction when the output is feasible, and by a wider margin for the without-reconstruction and count variants. For very simple or sparse cyclic queries no intermediate explosion occurs, so the optimiser's flat plan may already be good and the decomposition's overhead may dominate; we therefore expect the decomposition's advantage to grow with data density and scale.

    \item[\hypertarget{h5}{H5 (Semiring / FAQ aggregation).}] Because we expect the without-reconstruction decomposition to be fast whereas the final re-join of all bags to be expensive, we hypothesise that a semiring formulation --- mapping joins to multiplication and projections to addition over an implicit per-tuple count of $1$ --- can deliver the full result \emph{count} at close to the (cheap) without-reconstruction cost, avoiding the (expensive) materialisation of the full output. The same machinery generalises to other semirings (e.g.\ a $(\min,+)$ semiring for minimum-cost paths).
\end{description}

\section{Results and Discussion}
\label{sec:results_discussion}

We discuss the results by hypothesis, drawing on Tables~\ref{tab:results_soc-Epinions1}--\ref{tab:results_imdb_job}. 

%\subsection{H1 --- Tree Orientation}
%\label{sec:disc_h1}
%
%The orientation effect is large but its \emph{direction} is query-dependent, so \hyperlink{h1}{H1} holds only partially. The decisive factor is which bag becomes the heavy root and how large the intermediate it materialises is, rather than balancedness per se.
%
%The more balanced orientation \textbf{D2} wins where it produces a smaller or comparable root intermediate: on \texttt{cit-HepTh} Cyclic Query~4, $D_2'^{\text{full}}$ completes in $19.07$\,min against $24.08$\,min for $D_1'^{\text{full}}$, and the without-reconstruction $D_2'$ runs in $3.16$\,min versus $6.03$\,min for $D_1'$; on \texttt{soc-Epinions1} Cyclic Query~5, $D_2'$ ($11.34$\,min) is roughly twice as fast as $D_1'$ ($24.19$\,min).
%
%The less balanced orientation \textbf{D1} wins, however, when the balanced root happens to aggregate a much larger intermediate. On \texttt{soc-Epinions1} Cyclic Query~2 the root bag of $D_2'$ holds $28{,}048{,}050$ rows against only $3{,}544{,}194$ for $D_1'$, and the full count consequently takes $12.93$\,min for $D_2'^{\text{count}}$ versus $1.77$\,min for $D_1'^{\text{count}}$ --- a $7\times$ gap in favour of \textbf{D1}. This reflects a hard limit on what orientation can achieve: the $P'$ semijoins only remove dangling tuples and cannot shrink a bag below its \emph{semantic size} --- the number of distinct tuples its variables take in the true answer. When an orientation places a semantically large bag (here a dense two-atom ``wedge'' pattern) at the root, no amount of filtering can rescue it. A similar reversal appears on Cyclic Query~1 of the same dataset ($D_1'$ $0.85$\,min vs $D_2'$ $1.47$\,min). The hypothesis that a balanced orientation is \emph{uniformly} better is thus not supported: orientation should be chosen per query based on the resulting root intermediate size.
%
%\subsection{H2 --- Input-Connector Strategy}
%\label{sec:disc_h2}
%
%\hyperlink{h2}{H2} is confirmed once throughput inflation is accounted for. Across the graph datasets the wall-clock time of the single- and multi-connector variants is essentially identical: e.g.\ the 5-path query on \texttt{soc-Epinions1} completes in $0.10$\,min under both the single and the five-connector configuration, and on \texttt{p2p-Gnutella31} the baseline path query runs in $0.19$\,min in both cases. The throughput figures differ by roughly the connector count (e.g.\ $81{,}663$\,r/s vs $410{,}822$\,r/s), but this is an artefact of counting the re-ingested edge file once per connector (Section~\ref{sec:caveats}), not a genuine speed-up.
%
%Where a difference exists, the multi-connector variant is marginally \emph{slower} and uses somewhat more memory, consistent with the extra ingestion: on \texttt{cit-HepTh} Cyclic Query~1, $D_1^{\text{full}}$ takes $0.94$\,min ($7.7$\,GiB) with a single connector against $1.11$\,min ($9.5$\,GiB) with nine. This supports the hypothesis that processing, not ingestion, dominates, with the single-connector variant being the slightly leaner choice.
%
%\subsection{H3 --- Acyclic Queries}
%\label{sec:disc_h3}
%
%The acyclic results (TPC-H and IMDB) support \hyperlink{h3}{H3} on all counts.
%
%\paragraph{Materialisation overhead of full reconstruction.} On TPC-H the full-reconstruction decomposition matches the baseline in time but not in memory, and the gap widens with scale. At scale factor~10 (4-table query) the baseline runs in $10.34$\,min using $16.2$\,GiB, while $D_1^{\text{full}}$ takes $12.29$\,min and $42.2$\,GiB --- a $2.6\times$ memory overhead for no speed-up. At scale factor~30 both full-reconstruction variants ($D_1^{\text{full}}$ and $D_1'^{\text{full}}$) crash with out-of-memory, confirming that the forced materialisation of all bag views is the binding constraint.
%
%\paragraph{The lean variants match or beat the baseline.} The without-reconstruction, no-projection variant is consistently the most efficient. At scale factor~10 (4-table) $D_1'$ completes in $9.95$\,min using only $12.4$\,GiB --- faster and far leaner than the baseline --- and at scale factor~30 it finishes in $29.98$\,min where the full variants crash. This matches the expectation that, for $w(Q)=1$ queries, the decomposition's value is in bounding intermediates, not in reconstructing the full output.
%
%\paragraph{Semijoins add redundant work.} Comparing the projection and no-projection variants isolates the cost of the inline semijoins. On IMDB Acyclic Query~2 the no-projection variants are uniformly faster than their projection counterparts ($D_1'$ at $1.10$\,min vs $D_1$ at $1.18$\,min; $D_1'^{\text{full}}$ at $1.91$\,min vs $D_1^{\text{full}}$ at $2.02$\,min), consistent with the redundant-semijoin overhead discussed in Section~\ref{sec:decomp_no_semijoin}.
%
%\paragraph{Decomposition can beat a suboptimal baseline.} On IMDB Acyclic Query~2 the decomposition is markedly faster than the baseline ($D_1'$ $1.10$\,min vs baseline $2.91$\,min at comparable memory), indicating that the optimiser's flat-join plan is suboptimal and that the bag structure bounds intermediates more effectively. Conversely, the trivial IMDB Acyclic Query~1 shows no meaningful difference between any variants (all $\approx 0.18$\,min), as expected when the query is small enough that plan choice is irrelevant.
%
%\subsection{H4 --- Cyclic Queries}
%\label{sec:disc_h4}
%
%The cyclic results strongly support \hyperlink{h4}{H4}, including its caveat for simple/sparse instances.
%
%\paragraph{Dense graphs: large wins.} On the denser \texttt{soc-Epinions1} and \texttt{cit-HepTh} graphs the baseline frequently times out or runs orders of magnitude slower, while the decomposition completes quickly. On \texttt{cit-HepTh} Cyclic Query~2 the baseline takes $80.77$\,min against $1.88$\,min for $D_1'^{\text{full}}$ --- a $43\times$ speed-up at comparable memory --- and on \texttt{soc-Epinions1} Cyclic Query~1 the baseline times out entirely while the decomposition finishes in about a minute. This is the join-explosion-bounding behaviour predicted by the $w(Q)>1$ analysis.
%
%\paragraph{Sparse graphs: baseline competitive.} On the sparse \texttt{p2p-Gnutella31}, where the cyclic queries return very few rows (e.g.\ $0$, $1$, or $51$ result tuples), no join explosion occurs and the baseline is competitive or faster: on Cyclic Query~4 the baseline runs in $0.11$\,min against $0.35$\,min for $D_1'^{\text{full}}$, the decomposition's extra views now being pure overhead. This confirms the hypothesised regime in which the optimiser's plan is already good and the decomposition does not pay off.
%
%\subsection{H5 --- Semiring / FAQ Aggregation}
%\label{sec:disc_h5}
%
%\hyperlink{h5}{H5} is confirmed: the semiring count variant delivers the full result count at close to the cheap without-reconstruction cost, far below the cost of full materialisation. On \texttt{cit-HepTh} Cyclic Query~2, $D_1'^{\text{count}}$ returns the correct $266{,}404{,}327$-row count in $0.52$\,min using $7.7$\,GiB, against $1.88$\,min and $16.1$\,GiB for the full-reconstruction $D_1'^{\text{full}}$ --- roughly $3.6\times$ faster at half the memory --- while remaining close to the without-reconstruction $D_1'$ time of $0.44$\,min. The advantage is starker where full reconstruction is infeasible: on \texttt{soc-Epinions1} Cyclic Query~5, $D_1'^{\text{full}}$ times out, yet $D_1'^{\text{count}}$ completes in $41.83$\,min and returns the full count. This demonstrates that the fast, bounded decomposition can be extended with semirings to recover aggregate answers without paying the reconstruction cost.
%
%The generalisation to other semirings --- attaching a random integer cost in $[1,10]$ to every edge and computing, via a $(\min,+)$ semiring (joins $\mapsto +$, projections $\mapsto \min$), the minimum-cost 3-path per origin, summed for a single reported figure --- is in progress and will be reported alongside the count results. Likewise, the two additional TPC-H queries whose output cardinality differs from $|\texttt{LINEITEM}|$ will test whether the variant equivalence observed on the current TPC-H queries persists once the output decouples from the heaviest relation.
%
%\section{Summary of Findings and Threats to Validity}
%\label{sec:eval_summary}
%
%\paragraph{Summary.} The decomposition's benefit is concentrated on cyclic, dense queries (H4) and on acyclic queries where the optimiser's plan is poor (H3); on well-optimised acyclic and sparse cyclic queries it offers no advantage and pays a materialisation overhead. Full reconstruction is the main liability --- it inflates memory and can crash at scale --- whereas the without-reconstruction and semiring-count variants are consistently efficient (H3, H5). Orientation matters but must be chosen per query rather than by a fixed balancedness rule (H1), and the single-connector ingestion strategy is the marginally leaner choice (H2).
%
%\paragraph{Threats to validity.} The measurement caveats of Section~\ref{sec:caveats} (averaging noise, coarse memory sampling, connector-dependent throughput) bound the precision of small differences. Results are specific to Feldera's DBSP engine and its optimiser; a different engine might choose different baseline plans and shift the H3/H4 boundaries. Finally, several cells remain to be (re-)measured, and the additional TPC-H and IMDB queries are still in progress; the conclusions above are therefore preliminary where so noted.
%